\section{问题描述}

给定足够光滑且有界的关节期望轨迹 \(q_d(t),\dot q_d(t),\ddot q_d(t)\)。定义关节位置误差和速度误差为
\(
e=q_m-q_d\),\(\dot e=\dot q_m-\dot q_d\),
并令误差状态
\(
z=[e^\top\ \dot e^\top]^\top\in\mathbb R^{2n}.
\)
本文目标是在捕获前自由漂浮条件下设计关节力矩 \(\tau\)，使闭环误差在预设时间 \(T\) 后满足
\[
e(t)=0,\qquad \dot e(t)=0,\qquad t\ge T,
\]
或在存在扰动和模型近似时收敛到零点附近的小邻域。为保证 PDF 固定时间镇定结论具有明确的适用对象，本文不将延迟反馈直接叠加到原非线性耦合动力学上，而是通过模型补偿建立从辅助输入到误差状态的等效线性通道，并在该误差层完成周期延迟反馈设计。

本文采用如下假设。

\textbf{假设1：} \(M_e(q_m)\) 在工作空间内一致正定且可逆，\(M_e(q_m)\) 与 \(C_e(q_m,\dot q_m)\) 可由模型获得或可由名义模型近似。

\textbf{假设2：} 参考轨迹 \(q_d,\dot q_d,\ddot q_d\) 有界且连续，关节位置和速度可测。

\textbf{假设3：} 人工延迟项 \(z(t-h)\) 可由历史缓存获得，其中 \(h>0\) 为设计延迟。

\textbf{假设4：} 若采用预设时间扰动观测器，则等效加速度扰动 \(\Delta_a=M_e^{-1}(q_m)d(t)\) 及其导数有界，但上界可未知。

\textbf{假设5：} 对于严格固定时间精确归零结论，暂不考虑执行器饱和和模型误差，且扰动能够被精确补偿；若存在观测误差、模型近似或持续未补偿扰动，则讨论实际固定时间收敛性能。

上述假设将本文的理论结论限定在理想补偿和可获得历史误差的情形下。对于实际系统，传感器噪声、采样周期、执行器饱和和动力学建模误差会破坏严格线性闭环结构，因此仿真部分同时给出有界扰动下的实际跟踪结果。
